% Source: Wald GR, physical PDF pp. 192--193
%         (printed pp. 179--180).
% Coverage: Section 7.3, Algebraically Special Solutions; equation (7.3.1)
%           and Table 7.1.
% Figures/tables/footnotes: Table 7.1; no figures or footnotes.
% Status: source-reviewed and compile-clean.
% Uncertainties: none.

\subsection{Algebraically Special Solutions}\label{sec:7.3}

In the previous two sections (as well as in chapters 5 and 6) progress was made
toward obtaining physically interesting solutions of Einstein's equation by
restricting attention to spacetimes with a high degree of symmetry. In this section,
we shall very briefly discuss a different type of simplifying assumption for obtaining
solutions, which relates to the algebraic properties of the curvature tensor.
Unfortunately, since the character of this simplifying assumption is more
mathematical than physical, many of the solutions obtained by this approach do not
appear to be of direct physical relevance. Nevertheless, this approach has been one
of the most successful in providing us with a large class of exact solutions, and some
physically very important solutions, such as the Kerr metric (see chapter 12), have
been found in this way.

In the same manner as for the spatial metrics considered at the beginning of
section 5.1, at each point in spacetime we may view the Riemann tensor
\(R_{ab}{}^{cd}\) as a linear map from the six-dimensional space of antisymmetric
tensors of type \((0,2)\) (i.e., two-forms) into itself. This map is self-adjoint since
\(R_{abcd}=R_{cdab}\). However, it should be noted that unlike the case of a
Riemannian metric considered in section 5.1, the inner product induced on the
two-forms by the spacetime metric \(g_{ab}\) is not positive definite, so the familiar
theorem that a self-adjoint linear map has a complete orthonormal basis of
eigenvectors (which holds in any finite-dimensional vector space with a positive
definite inner product) does not apply here.

The analysis of the structure of the Riemann tensor as a linear map was first carried
out by Petrov (1954, 1969). We shall not present the details of the analysis of the
curvature tensor here but turn immediately to an important conclusion of the
analysis: In general, there exist precisely four distinct null directions (i.e., null
vectors \(k^a\), defined up to scaling \(k^a\longmapsto\lambda k^a\)) which
satisfy the relation
\begin{equation}
  k^bk^ck_{[e}C_{a]bc[d}k_{f]}=0 ,
  \tag{7.3.1}\label{eq:7.3.1}
\end{equation}
where \(C_{abcd}\) is the Weyl tensor, defined by equation~\eqref{eq:3.2.28}.
The null directions which satisfy equation~\eqref{eq:7.3.1} are called
\emph{principal null directions}. Thus, every tensor satisfying the algebraic
conditions satisfied by the Weyl tensor over a four-dimensional vector space with a
metric of Lorentz signature possesses, in general, four principal null directions. A
proof of this result by tensor methods requires a considerable amount of analysis.
However, a simple proof of this result can be obtained by spinor methods, and we
will present this proof in chapter 13.

Although, in general, the Weyl tensor possesses four distinct principal null
directions, it is possible for some of these null directions to coincide---in which
case they satisfy stronger relations than equation~\eqref{eq:7.3.1}---resulting in
fewer than four principal null directions. Spacetimes for which fewer than four
distinct principal null directions exist at each point are called \emph{algebraically
special spacetimes}. The different types of algebraically special spacetimes are
classified in Table~\ref{tab:7.1}.

\begingroup
\renewcommand{\thetable}{7.1}
\begin{table}[htbp]
  \centering
  \caption{Algebraic classification of spacetimes}
  \label{tab:7.1}
  \small
  \renewcommand{\arraystretch}{1.3}
  \begin{tabular}{@{}c >{\centering\arraybackslash}p{0.35\linewidth}
                    >{\centering\arraybackslash}p{0.43\linewidth}@{}}
    \toprule
    Type & Description &
      Condition satisfied by (repeated) principal null direction \(k^a\) \\
    \midrule
    I &
      Algebraically general; four distinct principal null directions &
      \(k^bk^ck_{[e}C_{a]bc[d}k_{f]}=0\) \\
    II &
      One pair of principal null directions coincides &
      \(k^bk^cC_{abc[d}k_{e]}=0\) \\
    II--II [D] &
      Two pairs of principal null directions coincide &
      \(\begin{gathered}k^bk^cC_{abc[d}k_{e]}=0\\
        \text{(two solutions)}\end{gathered}\) \\
    III &
      Three principal null directions coincide &
      \(k^cC_{abc[d}k_{e]}=0\) \\
    IV [N] &
      All four principal null directions coincide &
      \(k^cC_{abcd}=0\) \\
    \bottomrule
  \end{tabular}
\end{table}
\endgroup

In algebraically special spacetimes, we can take advantage of the conditions
satisfied by the repeated principal null direction by choosing a null tetrad (see the
discussion of the Newman--Penrose formalism at the end of section 3.4) with one of
the null vectors aligned with the repeated principal null direction. Certain tetrad
components of the Weyl tensor then will vanish. This yields extra equations on the
connection components (i.e., the spin coefficients in the Newman--Penrose
formalism) in addition to the conditions imposed by Einstein's equation. Using
these additional conditions, it has proven possible to integrate Einstein's equation
exactly in many cases.

We shall not attempt to present here any of the details of how algebraically special
solutions can be found. For an excellent illustration of this approach, we refer the
reader to the paper of Kinnersley (1969) who explicitly found all type II--II solutions
of the vacuum Einstein equation by direct integration of the Newman--Penrose
equations. A survey of the algebraically special solutions known by the late 1970s
can be found in the book of Kramer et al.\ (1980).
